% 4.8 诱导表示的对称化与反对称化平方
\section{诱导表示的对称化与反对称化平方}

在 2.6 节我们引入了表示 $\Lambda$ 的对称化幂与反对称化幂的概念，并以晶体学点群表示的对称化与反对称化平方为例说明了这一普遍理论。本节将说明如何把诱导表示的 Kronecker 平方的空间分解为其对称部分与反对称部分。然后通过推广上一节的例子，把这一理论应用于空间群表示以演示其结果。金刚石结构（$Fd3m$, $O^{7}_{h}$）与闪锌矿结构（$F\bar{4}3m$, $T^{2}_{d}$）的布里渊区某些点的空间群表示的对称化平方与立方，已由 Birman (1962\textit{b}) 列表；六角密堆积（$P6_{3}/mmc$, $D^{4}_{6h}$）与纤锌矿（$P6_{3}mc$, $C^{4}_{6v}$）结构的对称化平方已由 Chen and Hsieh (1965) 列表。本节提出的理论是 Mackey (1952) 的工作的推广，不过视角略有不同（Bradley and Davies 1970）；特别地，我们的证明是构造性的，并给出执行分解的一个明确处方。

我们简略地说明对称化积与反对称化积的概念在量子力学选择定则研究中的物理应用。若设系统哈密顿量的幺正对称操作群是空间群 $\mathbf{G}$，那么我们在上一节已经看到如何确定系统中可能发生的各种过程的选择定则。如果现在把时间反演对称操作 $\theta$ 加进来以扩充群 $\mathbf{G}$，那么 4.7 节通常的三重 Kronecker 积中初态与末态的乘积 $\Gamma^{i} \boxtimes \Gamma^{k}$ 可能不得不换成对称化或反对称化乘积。这发生在末态与初态的时间反演相联系的时候，其效果是可能提供额外的选择定则。证明这一点的原始论文是关于 Jahn--Teller 效应的（Jahn 1938, Jahn and Teller 1937），但下述由 Lax (1962) 给出的一小段理论看来已经足够。若末态的波函数 $\Psi^{i}_{r}$ 通过时间反演对称与初态相联系，即 $\Psi^{i}_{r} = \theta\Psi^{k}_{r}$，则考虑矩阵元
\begin{equation*}
W^{k}_{rt} = \langle \theta\Psi^{k}_{r}, W\Psi^{k}_{t}\rangle
\tag{4.8.1}
\end{equation*}
\begin{equation*}
= \langle \theta W\Psi^{k}_{t}, \theta^{2}\Psi^{k}_{r}\rangle
\tag{4.8.2}
\end{equation*}
因为 $\theta$ 是反幺正的（见第七章）。现在 $\Delta(\theta^{2}) = \omega = \pm 1$，并作一个限制性不强的假定 $(\theta W\theta^{-1})^{\dagger} = \alpha W$，其中剑号表示伴随算符，$\alpha = \pm 1$。把这些关系代入式 (4.8.2) 得
\begin{equation*}
W^{k}_{rt} = \alpha\omega W^{k}_{tr} = \pm W^{k}_{tr}.
\tag{4.8.3}
\end{equation*}
由于 $\Psi^{k}_{t}$ 与 $\Phi^{k}_{r} = (\theta\Psi^{k}_{r})^{*}$ 都属于 $\Gamma^{k}$，由式 (4.8.3) 可得结论：此情形下选择定则的存在与否，取决于积分 $\int(\Phi^{k}_{r}W\Psi^{k}_{t} \pm \Phi^{k}_{t}W\Psi^{k}_{r})\,\mathrm{d}x$ 在 $\mathbf{G}$ 的操作下的行为。函数 $(\Phi^{k}_{r}\Psi^{k}_{t} \pm \Phi^{k}_{t}\Psi^{k}_{r})$ 张成的子空间的不变性，加上通常的选择定则理论，表明：当式 (4.8.3) 中取正号时，若乘积 $[\Gamma^{k} \boxtimes \Gamma^{k}] \boxtimes \Gamma^{w}$（取负号时，若乘积 $\{\Gamma^{k} \boxtimes \Gamma^{k}\} \boxtimes \Gamma^{w}$）不含 $\mathbf{G}$ 的恒等表示，则该跃迁被对称性禁戒。一个等价的判据是：$\Gamma^{k}$ 的对称化（或反对称化）平方是否含有 $\Gamma^{w*}$ 的某些不可约分量。上述理论可以从 $\Psi^{i}_{r} = \theta\Psi^{k}_{r}$ 的情形推广到 $\Psi^{i}_{r} = \{R \mid \mathbf{v}\}\theta\Psi^{k}_{r}$ 的情形，其中 $\{R \mid \mathbf{v}\}$ 是空间群 $\mathbf{G}$ 的任一元素（见 Lax (1962)）。另一个需要用给定表示的对称化幂与反对称化幂来分析的物理应用，是二级相变的 Landau 理论，2.6 节开头已提到（文献见 2.6 节）。

现在研究群 $\mathbf{G}$ 的诱导表示的对称化与反对称化平方。也就是说，要把 2.6 节记号中的 $\mathbf{S}$ 认同为群 $\mathbf{G}$，$\mathbf{H}$ 为 $\mathbf{G}$ 的子群，$\Lambda$（2.6 节记号）认同为由 $\mathbf{H}$ 的表示 $\mathbf{D}$ 诱导的 $\mathbf{G}$ 的表示。$\mathbf{D} \uparrow \mathbf{G}$ 的 Kronecker 平方的分解由上一节的理论覆盖，主要结果是式 (4.7.25)：
\begin{equation*}
(\mathbf{D} \uparrow \mathbf{G}) \boxtimes (\mathbf{D} \uparrow \mathbf{G}) \equiv \sum_{\alpha}\left[(D_{\alpha} \boxtimes D) \downarrow \mathbf{L}_{\alpha}\right] \uparrow \mathbf{G}.
\tag{4.8.4}
\end{equation*}
这里对 $\alpha$ 求和的各项与双陪集分解
\begin{equation*}
\mathbf{G} = \sum_{\alpha}\mathbf{H}d_{\alpha}\mathbf{H}
\tag{4.8.5}
\end{equation*}
一一对应，且 $\mathbf{L}_{\alpha} = \mathbf{H} \cap \mathbf{H}_{\alpha}$，其中 $\mathbf{H}_{\alpha} = d_{\alpha}\mathbf{H}d^{-1}_{\alpha}$。同时 $D_{\alpha}$ 是 $\mathbf{H}_{\alpha}$ 的表示，满足 $\mathbf{D}_{\alpha}(d_{\alpha}hd^{-1}_{\alpha}) = \mathbf{D}(h)$ 对一切 $h \in \mathbf{H}$。我们需要重新安排或修改式 (4.8.4) 的右端，使得其中一部分可以等同于对称化平方 $\left[(\mathbf{D} \uparrow \mathbf{G}) \boxtimes (\mathbf{D} \uparrow \mathbf{G})\right]$，另一部分等同于反对称化平方 $\{(\mathbf{D} \uparrow \mathbf{G}) \boxtimes (\mathbf{D} \uparrow \mathbf{G})\}$。我们沿用上一节的记号：$\mathbf{D}$ 是 $\mathbf{H}$ 的维数为 $d$ 的表示，基 $\langle \phi | = \langle \phi_{1}, \phi_{2}, \dots, \phi_{d} |$，而 $D_{\alpha}$ 是 $\mathbf{H}_{\alpha}$ 的表示，基为 $d_{\alpha}\langle \phi | = \langle \phi_{\alpha} | = \langle \phi_{\alpha 1}, \phi_{\alpha 2}, \dots, \phi_{\alpha d} |$。确实
\begin{equation*}
\begin{array}{rl}
d_{\alpha}hd^{-1}_{\alpha}\langle \phi_{\alpha} | &= d_{\alpha}h\langle \phi |\\
&= d_{\alpha}\langle \phi |\mathbf{D}(h)\\
&= \langle \phi_{\alpha} |\mathbf{D}(h),
\end{array}
\tag{4.8.6}
\end{equation*}
从而验证了关系 $\mathbf{D}_{\alpha}(d_{\alpha}hd^{-1}_{\alpha}) = \mathbf{D}(h)$。表示 $(\mathbf{D} \uparrow \mathbf{G})$ 的基是 $p_{\sigma}\langle \phi |$（$\sigma = 1$ 到 $|\mathbf{G}|/|\mathbf{H}|$），其中 $p_{\sigma}$ 是分解 $\mathbf{G} = \sum_{\sigma}p_{\sigma}\mathbf{H}$ 中的陪集代表。另外，由上一节记得陪集分解 $\mathbf{H} = \sum_{\gamma}q_{\alpha\gamma}\mathbf{L}_{\alpha}$，于是固定 $\alpha$ 时 $\left[(D_{\alpha} \boxtimes D) \downarrow \mathbf{L}_{\alpha}\right] \uparrow \mathbf{G}$ 的基是函数集 $p_{\sigma}q_{\alpha\gamma}(\phi_{\alpha i}, \phi_{j})$（$i = 1$ 到 $d$，$j = 1$ 到 $d$，$\gamma = 1$ 到 $|\mathbf{H}|/|\mathbf{L}_{\alpha}|$，$\sigma = 1$ 到 $|\mathbf{G}|/|\mathbf{H}|$）。依据定理 4.7.7，由这 $d^{2}|\mathbf{G}|/|\mathbf{L}_{\alpha}|$ 个函数张成的矢量空间 $V_{\alpha}$ 在 $(\mathbf{D} \uparrow \mathbf{G}) \boxtimes (\mathbf{D} \uparrow \mathbf{G})$ 下不变。可以符号化地写
\begin{equation*}
V_{\alpha} = \sum_{\sigma, \gamma, i, j}p_{\sigma}q_{\alpha\gamma}(\phi_{\alpha i}, \phi_{j}),
\tag{4.8.7}
\end{equation*}
其中右端指函数 $p_{\sigma}q_{\alpha\gamma}(\phi_{\alpha i}, \phi_{j})$ 的所有线性组合的集合。空间 $V_{\alpha}$ 与所选定的双陪集代表无关：如果不用 $d_{\alpha}$ 而用 $d_{\beta} = h_{a}d_{\alpha}h_{b}$（$h_{a}, h_{b} \in \mathbf{H}$），则 $V_{\beta}$ 与 $V_{\alpha}$ 重合。这一点在 4.7 节是隐含的，可以如下看出：
\begin{equation*}
\begin{array}{rl}
\mathbf{L}_{\beta} &= \mathbf{H} \cap \mathbf{H}_{\beta} = \mathbf{H} \cap d_{\beta}\mathbf{H}d^{-1}_{\beta} = \mathbf{H} \cap h_{a}d_{\alpha}h_{b}\mathbf{H}h^{-1}_{b}d^{-1}_{\alpha}h^{-1}_{a}\\
&= \mathbf{H} \cap h_{a}d_{\alpha}\mathbf{H}d^{-1}_{\alpha}h^{-1}_{a} = \mathbf{H} \cap h_{a}\mathbf{H}_{\alpha}h^{-1}_{a},
\end{array}
\end{equation*}
因此 $h^{-1}_{a}\mathbf{L}_{\beta}h_{a} = h^{-1}_{a}\mathbf{H}h_{a} \cap \mathbf{H}_{\alpha} = \mathbf{H} \cap \mathbf{H}_{\alpha} = \mathbf{L}_{\alpha}$。于是若 $\mathbf{H} = \sum_{\gamma}q_{\alpha\gamma}\mathbf{L}_{\alpha}$，则也有 $\mathbf{H}h^{-1}_{a} = \sum_{\gamma}q_{\alpha\gamma}h^{-1}_{a}\mathbf{L}_{\beta}$，从而对每个 $\gamma$ 可取 $q_{\beta\gamma} = q_{\alpha\gamma}h^{-1}_{a}$，并且
\begin{equation*}
\begin{array}{rl}
V_{\beta} &= \sum_{\sigma, \gamma, i, j}p_{\sigma}q_{\beta\gamma}(\phi_{\beta i}, \phi_{j})\\
&= \sum_{\sigma, \gamma, i, j}p_{\sigma}q_{\alpha\gamma}h^{-1}_{a}(d_{\beta}\phi_{i}, \phi_{j})\\
&= \sum_{\sigma, \gamma, i, j}p_{\sigma}q_{\alpha\gamma}(d_{\alpha}h_{b}\phi_{i}, h^{-1}_{a}\phi_{j})\\
&= \sum_{\sigma, \gamma, k, l}p_{\sigma}q_{\alpha\gamma}(\phi_{\alpha k}, \phi_{l})\\
&= V_{\alpha}.
\end{array}
\tag{4.8.8}
\end{equation*}
矢量空间 $V$ 是整个表示 $(\mathbf{D} \uparrow \mathbf{G}) \boxtimes (\mathbf{D} \uparrow \mathbf{G})$ 的承载空间，如式 (4.8.5) 那样对不同的双陪集求和，它当然已被证明是直和 $\sum_{\alpha}V_{\alpha}$。因此除非 $d_{\alpha}$ 与 $d_{\beta}$ 属于同一双陪集，任何两个矢量空间 $V_{\alpha}$ 与 $V_{\beta}$ 都不会重合。

由于即将显见的原因，给定一个特定的双陪集代表 $d_{\alpha}$，定义 $d^{-1}_{\alpha}\phi_{i} \equiv \phi_{\bar{\alpha}i}$ 并写 $\mathbf{L}_{\bar{\alpha}} = \mathbf{H} \cap \mathbf{H}_{\bar{\alpha}} = \mathbf{H} \cap d^{-1}_{\alpha}\mathbf{H}d_{\alpha}$ 是有用的。注意 $\mathbf{L}_{\bar{\alpha}} = d^{-1}_{\alpha}\mathbf{L}_{\alpha}d_{\alpha}$，并且由于 $\mathbf{G} = \sum_{\sigma}p_{\sigma}\mathbf{H}$ 与 $\mathbf{H} = \sum_{\gamma}q_{\alpha\gamma}\mathbf{L}_{\alpha}$，有 $\mathbf{G} = \mathbf{G}d_{\alpha} = \sum_{\sigma, \gamma}p_{\sigma}q_{\alpha\gamma}\mathbf{L}_{\alpha}d_{\alpha} = \sum_{\sigma, \gamma}p_{\sigma}q_{\alpha\gamma}d_{\alpha}\mathbf{L}_{\bar{\alpha}}$。于是当 $\sigma$ 与 $\gamma$ 取遍所有可能值时，集合 $p_{\sigma}q_{\alpha\gamma}d_{\alpha}$ 构成 $\mathbf{L}_{\bar{\alpha}}$ 在 $\mathbf{G}$ 中的陪集代表的完全集。因此若把矢量空间 $W_{\bar{\alpha}}$ 定义为 $\left[(D \boxtimes D_{\alpha}) \downarrow \mathbf{L}_{\bar{\alpha}}\right] \uparrow \mathbf{G}$ 的承载空间，就会发现
\begin{equation*}
\begin{array}{rl}
W_{\bar{\alpha}} &= \sum_{\sigma, \gamma, i, j}p_{\sigma}q_{\alpha\gamma}d_{\alpha}(\phi_{i}, \phi_{\alpha j})\\
&= \sum_{\sigma, \gamma, i, j}p_{\sigma}q_{\alpha\gamma}(\phi_{\alpha i}, \phi_{j}) = V_{\alpha}.
\end{array}
\tag{4.8.9}
\end{equation*}
故矢量空间 $W_{\bar{\alpha}}$ 与 $V_{\alpha}$ 重合。类似地 $W_{\alpha}$ 与 $V_{\bar{\alpha}}$ 重合。现在 $\left[(D \boxtimes D_{\alpha}) \downarrow \mathbf{L}_{\bar{\alpha}}\right] \uparrow \mathbf{G}$ 的特征标显然与 $\left[(D_{\alpha} \boxtimes D) \downarrow \mathbf{L}_{\alpha}\right] \uparrow \mathbf{G}$ 的相同；由于前者定义在 $W_{\alpha} = V_{\bar{\alpha}}$ 上、后者定义在 $V_{\alpha}$ 上，所以分别定义在 $V_{\alpha}$ 与 $V_{\bar{\alpha}}$ 上的两个表示等价。现在可能发生两种情形：要么 $V_{\alpha}$ 与 $V_{\bar{\alpha}}$ 重合（于是通过平凡的等价变换定义等价表示），要么 $V_{\alpha}$ 与 $V_{\bar{\alpha}}$ 不同（但仍然定义等价表示）。由前面的讨论我们知道，$V_{\alpha}$ 与 $V_{\bar{\alpha}}$ 重合当且仅当 $\mathbf{H}d_{\alpha}\mathbf{H} = \mathbf{H}d^{-1}_{\alpha}\mathbf{H}$，因此发生哪种情形，关键取决于双陪集 $\mathbf{H}d_{\alpha}\mathbf{H}$ 是否\textit{自逆}（因为 $(\mathbf{H}d_{\alpha}\mathbf{H})^{-1} = \mathbf{H}d^{-1}_{\alpha}\mathbf{H}$，而 $\mathbf{H} = \mathbf{H}^{-1}$）。

先考虑 $V_{\alpha}$ 与 $V_{\bar{\alpha}}$ 不同的情形。现在 $V_{\alpha} = \sum_{\sigma, \gamma, i, j}p_{\sigma}q_{\alpha\gamma}(\phi_{\alpha i}, \phi_{j})$ 而 $V_{\bar{\alpha}} = W_{\alpha} = \sum_{\sigma, \gamma, i, j}p_{\sigma}q_{\alpha\gamma}(\phi_{i}, \phi_{\alpha j})$。它们是不同的空间，却定义等价的表示。显然可以作直和 $(V_{\alpha} \oplus V_{\bar{\alpha}})$，并把它分解成另一种直和 $(V^{+}_{\alpha} \oplus V^{-}_{\alpha})$，其中
\begin{equation*}
V^{\pm}_{\alpha} = \sum_{\sigma, \gamma, i, j}p_{\sigma}q_{\alpha\gamma}\left\{(\phi_{\alpha i}, \phi_{j}) \pm (\phi_{j}, \phi_{\alpha i})\right\}
\tag{4.8.10}
\end{equation*}
上下号分别对应。由式 (4.8.10) 右端的形式也清楚可见：$V^{+}_{\alpha}$ 是 $V^{+}$ 的子空间，$V^{-}_{\alpha}$ 是 $V^{-}$ 的子空间，其中 $V^{+}$ 与 $V^{-}$ 是 $V$ 本身的对称化与反对称化子空间。而且定义在 $V^{+}_{\alpha}$、$V^{-}_{\alpha}$、$V_{\alpha}$、$V_{\bar{\alpha}}$ 上的表示全都等价，前两个是由后两个通过式 (4.8.10) 表达的简单等价变换得到的。于是，若 $V_{\alpha}$ 与 $V_{\bar{\alpha}}$ 不同，即 $d_{\alpha}$ 与 $d^{-1}_{\alpha}$ 属于不同的双陪集，则 $[(\mathbf{D} \uparrow \mathbf{G}) \boxtimes (\mathbf{D} \uparrow \mathbf{G})]$ 与 $\{(\mathbf{D} \uparrow \mathbf{G}) \boxtimes (\mathbf{D} \uparrow \mathbf{G})\}$ 都将含有一个等价于 $\left[(D_{\alpha} \boxtimes D) \downarrow \mathbf{L}_{\alpha}\right] \uparrow \mathbf{G}$ 的表示。

再考虑 $V_{\alpha} = V_{\bar{\alpha}}$ 的情形。有两种可能需要分别处理，因为分解取决于 $D_{\alpha}$ 与 $D$ 的承载空间相同还是不同。这两个承载空间相同当且仅当 $d_{\alpha} \in \mathbf{H}$，此时可取 $d_{\alpha}$ 为恒等元 $d_{\varepsilon}$。于是 $\mathbf{H}_{\varepsilon} = \mathbf{H}$，$\mathbf{L}_{\varepsilon} = \mathbf{H} \cap \mathbf{H}_{\varepsilon} = \mathbf{H}$。式 (4.8.4) 中现在考虑的这一项就是 $(D \boxtimes D) \downarrow \mathbf{H}$ 诱导的 $(\mathbf{D} \boxtimes \mathbf{D}) \uparrow \mathbf{G}$，而 $V_{\varepsilon} = \sum_{\sigma, i, j}p_{\sigma}(\phi_{i}, \phi_{j})$。立即看出 $V_{\varepsilon}$ 分解为直和 $(V^{+}_{\varepsilon} \oplus V^{-}_{\varepsilon})$，其中
\begin{equation*}
V^{\pm}_{\varepsilon} = \sum_{\sigma, i, j}p_{\sigma}\left\{(\phi_{i}, \phi_{j}) \pm (\phi_{j}, \phi_{i})\right\}
\tag{4.8.11}
\end{equation*}
且 $V^{+}_{\varepsilon}$ 是 $V^{+}$ 的子空间、$V^{-}_{\varepsilon}$ 是 $V^{-}$ 的子空间。由式 (4.8.11) 可见，$V^{+}_{\varepsilon}$ 定义表示 $[\mathbf{D} \boxtimes \mathbf{D}] \uparrow \mathbf{G}$，$V^{-}_{\varepsilon}$ 定义表示 $\{\mathbf{D} \boxtimes \mathbf{D}\} \uparrow \mathbf{G}$，这两个表示的维数分别为 $\frac{1}{2}d(d+1)k$ 与 $\frac{1}{2}d(d-1)k$，其中 $k = |\mathbf{G}|/|\mathbf{H}|$。这两个表示的维数之差等于 $dk$。由于 $[(\mathbf{D} \uparrow \mathbf{G}) \boxtimes (\mathbf{D} \uparrow \mathbf{G})]$ 与 $\{(\mathbf{D} \uparrow \mathbf{G}) \boxtimes (\mathbf{D} \uparrow \mathbf{G})\}$ 的维数分别为 $\frac{1}{2}dk(dk+1)$ 与 $\frac{1}{2}dk(dk-1)$，其维数之差也等于 $dk$，可见这个差完全由 $V_{\varepsilon}$ 的分解所解释。这使得任何其他满足 $V_{\varepsilon} = V_{\alpha}$ 的空间 $V_{\alpha}$ 在 $V$ 的最终分解中以某种方式给出同样维数的空间变得合理。现在我们着手进行一些相当细致的分析来展示这实际上如何发生。

现在考虑 $\mathbf{H}d_{\alpha}\mathbf{H} = \mathbf{H}d^{-1}_{\alpha}\mathbf{H}$ 而 $d_{\alpha}$ 不能取为恒等元 $d_{\varepsilon}$ 的情形。因此存在 $h_{i}, h_{j}, h_{k}, h_{l} \in \mathbf{H}$ 使 $h_{i}d_{\alpha}h_{j} = h_{k}d^{-1}_{\alpha}h_{l}$，即 $d_{\alpha}h_{j}h^{-1}_{l} = h^{-1}_{i}h_{k}d^{-1}_{\alpha}$。于是 $d_{\alpha}\mathbf{H} \cap \mathbf{H}d^{-1}_{\alpha}$ 非空。设 $z$ 是这个交集的任一元素；则 $z$ 不在 $\mathbf{H}$ 中，且可写 $z = d_{\alpha}h = h'd^{-1}_{\alpha}$，$h, h' \in \mathbf{H}$。现在
\begin{equation*}
z^{2} = d_{\alpha}hh'd^{-1}_{\alpha} = h'h \in d_{\alpha}\mathbf{H}d^{-1}_{\alpha} \cap \mathbf{H} = \mathbf{L}_{\alpha}
\end{equation*}
并且进一步
\begin{equation*}
\begin{array}{l}
z\mathbf{H}z^{-1} = d_{\alpha}h\mathbf{H}h^{-1}d^{-1}_{\alpha} = \mathbf{H}_{\alpha},\\
z\mathbf{H}_{\alpha}z^{-1} = h'd^{-1}_{\alpha}d_{\alpha}\mathbf{H}d^{-1}_{\alpha}d_{\alpha}h'^{-1} = \mathbf{H}.
\end{array}
\end{equation*}
于是可以构造群
\begin{equation*}
\mathbf{M}_{\alpha} = \mathbf{L}_{\alpha} \oplus z\mathbf{L}_{\alpha}
\tag{4.8.12}
\end{equation*}
且 $\mathbf{L}_{\alpha}$ 是 $\mathbf{M}_{\alpha}$ 的指数为 2 的不变子群。从 $\mathbf{L}_{\alpha}$ 上带基 $\langle \phi_{\alpha}|\langle \phi |$ 的表示 $(D_{\alpha} \boxtimes D)$ 出发，可以构造诱导表示 $C_{\alpha} = (D_{\alpha} \boxtimes D) \uparrow \mathbf{M}_{\alpha}$，其基为 $(\langle \phi_{\alpha}|\langle \phi |,\ z\langle \phi_{\alpha}|\langle \phi |)$。分析类似于 4.5 节 $p = 2$、$r = z$ 的情形，故略去细节，只引用结果。若 $l \in \mathbf{L}_{\alpha}$，则诱导表示由块矩阵表达式给出：
\begin{equation*}
\begin{array}{rl}
\mathbf{C}_{\alpha}(l) &= \begin{pmatrix} \mathbf{D}(d^{-1}_{\alpha}ld_{\alpha}) \otimes \mathbf{D}(l) & \mathbf{0}\\ \mathbf{0} & \mathbf{D}(d^{-1}_{\alpha}z^{-1}lzd_{\alpha}) \otimes \mathbf{D}(z^{-1}lz) \end{pmatrix},\\
\mathbf{C}_{\alpha}(z) &= \begin{pmatrix} \mathbf{0} & \mathbf{D}(d^{-1}_{\alpha}z^{2}d_{\alpha}) \otimes \mathbf{D}(z^{2})\\ \mathbf{1} & \mathbf{0} \end{pmatrix}.
\end{array}
\tag{4.8.13}
\end{equation*}
为方便起见，记
\begin{equation*}
\mathbf{N}_{\alpha}(l) = \mathbf{D}(d^{-1}_{\alpha}ld_{\alpha}) \otimes \mathbf{D}(l)
\tag{4.8.14}
\end{equation*}
以及
\begin{equation*}
{}_{z}\mathbf{N}_{\alpha}(l) = \mathbf{N}_{\alpha}(z^{-1}lz),
\tag{4.8.15}
\end{equation*}
使得
\begin{equation*}
\mathbf{C}_{\alpha}(l) = \begin{pmatrix} \mathbf{N}_{\alpha}(l) & \mathbf{0}\\ \mathbf{0} & {}_{z}\mathbf{N}_{\alpha}(l) \end{pmatrix},
\tag{4.8.16}
\end{equation*}
\begin{equation*}
\mathbf{C}_{\alpha}(z) = \begin{pmatrix} \mathbf{0} & \mathbf{N}_{\alpha}(z^{2})\\ \mathbf{1} & \mathbf{0} \end{pmatrix}.
\tag{4.8.17}
\end{equation*}
由 4.5 节的分析可知，如果能够找到一个矩阵 $\mathbf{P}$ 使
\begin{equation*}
\mathbf{P}^{2} = \mathbf{N}_{\alpha}(z^{2})
\tag{4.8.18}
\end{equation*}
且
\begin{equation*}
\mathbf{P}\,{}_{z}\mathbf{N}_{\alpha}(l)\mathbf{P}^{-1} = \mathbf{N}_{\alpha}(l)
\tag{4.8.19}
\end{equation*}
对一切 $l \in \mathbf{L}_{\alpha}$ 成立，则 ${}_{z}\mathbf{N}_{\alpha}$ 与 $\mathbf{N}_{\alpha}$ 等价；并且相对于变换后的基 $(\langle \phi_{\alpha}|\langle \phi |,\ z\langle \phi_{\alpha}|\langle \phi |\mathbf{P}^{-1})$，表示 $C_{\alpha}$ 变换为等价表示 $C'_{\alpha}$：
\begin{equation*}
\mathbf{C}'_{\alpha}(l) = \begin{pmatrix} \mathbf{N}_{\alpha}(l) & \mathbf{0}\\ \mathbf{0} & \mathbf{N}_{\alpha}(l) \end{pmatrix},
\tag{4.8.20}
\end{equation*}
\begin{equation*}
\mathbf{C}'_{\alpha}(z) = \begin{pmatrix} \mathbf{0} & \mathbf{P}\\ \mathbf{P} & \mathbf{0} \end{pmatrix}.
\tag{4.8.21}
\end{equation*}
满足式 (4.8.18) 与 (4.8.19) 的矩阵 $\mathbf{P}$ 可以找到，由下式给出：
\begin{equation*}
\mathbf{P}_{ts, uv} = \mathbf{D}(h')_{su}\mathbf{D}(h)_{tv} = \left\{\mathbf{D}(h') \otimes \mathbf{D}(h)\right\}_{st, uv}.
\tag{4.8.22}
\end{equation*}
例如，用这一选取可以很快验证式 (4.8.18) 成立：
\begin{equation*}
\begin{array}{rl}
\mathbf{P}^{2}_{ab, cd} &= \sum_{e, f}\mathbf{P}_{ab, ef}\mathbf{P}_{ef, cd}\\
&= \sum_{e, f}\mathbf{D}(h')_{be}\mathbf{D}(h)_{af}\mathbf{D}(h')_{fc}\mathbf{D}(h)_{ed}\\
&= \mathbf{D}(h'h)_{bd}\mathbf{D}(hh')_{ac}\\
&= \mathbf{D}(z^{2})_{bd}\mathbf{D}(d^{-1}_{\alpha}z^{2}d_{\alpha})_{ac}\\
&= \left\{\mathbf{D}(d^{-1}_{\alpha}z^{2}d_{\alpha}) \otimes \mathbf{D}(z^{2})\right\}_{ab, cd}\\
&= \mathbf{N}_{\alpha}(z^{2})_{ab, cd}.
\end{array}
\end{equation*}
为验证式 (4.8.19)，我们计算 $\mathbf{P}\,{}_{z}\mathbf{N}_{\alpha}(l)$ 并证明它等于 $\mathbf{N}_{\alpha}(l)\mathbf{P}$：
\begin{equation*}
\begin{array}{l}
\left(\mathbf{P}\,{}_{z}\mathbf{N}_{\alpha}(l)\right)_{ab, cd} = \sum_{e, f}\mathbf{P}_{ab, ef}\,{}_{z}\mathbf{N}_{\alpha}(l)_{ef, cd}\\
= \sum_{e, f}\mathbf{D}(h')_{be}\mathbf{D}(h)_{af}\mathbf{D}(d^{-1}_{\alpha}z^{-1}lzd_{\alpha})_{ec}\mathbf{D}(z^{-1}lz)_{fd}\\
= \mathbf{D}(lzd_{\alpha})_{bc}\mathbf{D}(d^{-1}_{\alpha}lz)_{ad}\\
= \sum_{e, f}\mathbf{D}(l)_{bf}\mathbf{D}(h')_{fc}\mathbf{D}(d^{-1}_{\alpha}ld_{\alpha})_{ae}\mathbf{D}(h)_{ed}\\
= \sum_{e, f}\mathbf{N}_{\alpha}(l)_{ab, ef}\mathbf{P}_{ef, cd}\\
= \left(\mathbf{N}_{\alpha}(l)\mathbf{P}\right)_{ab, cd}.
\end{array}
\tag{4.8.23}
\end{equation*}
在推导这些等式时我们用了表达式 $z = d_{\alpha}h = h'd^{-1}_{\alpha}$ 以及 $\mathbf{D}$ 是 $\mathbf{H}$ 的元素上的同态这一事实。若不是由于 $C'_{\alpha}$ 的基现在取得了一个特别简单的形式，这一切分析都将毫无用处。为看清这一点，计算基的第二个成员的 $(kl)$ 元：
\begin{equation*}
\begin{array}{rl}
\left(z\langle \phi_{\alpha}|\langle \phi |\mathbf{P}^{-1}\right)_{kl} &= \sum_{i, j}z(\phi_{\alpha i}, \phi_{j})\mathbf{P}^{-1}_{ij, kl}\\
&= \sum_{i, j}(h'd^{-1}_{\alpha}\phi_{\alpha i}, d_{\alpha}h\phi_{j})\mathbf{P}^{-1}_{ij, kl}\\
&= \sum_{i, j, t}(h'\phi_{i}, d_{\alpha}\phi_{t}\mathbf{D}(h)_{ij})\mathbf{P}^{-1}_{ij, kl}\\
&= \sum_{i, j, t, s}(\phi_{s}, \phi_{\alpha t})\mathbf{D}(h')_{si}\mathbf{D}(h)_{tj}\mathbf{P}^{-1}_{ij, kl}\\
&= \sum_{i, j, t, s}(\phi_{s}, \phi_{\alpha t})\mathbf{P}_{ts, ij}\mathbf{P}^{-1}_{ij, kl}\\
&= (\phi_{l}, \phi_{\alpha k}).
\end{array}
\tag{4.8.24}
\end{equation*}
另外
\begin{equation*}
\left(\langle \phi_{\alpha}|\langle \phi |\right)_{kl} = (\phi_{\alpha k}, \phi_{l}).
\tag{4.8.25}
\end{equation*}
于是
\begin{equation*}
\left(\langle \phi_{\alpha}|\langle \phi | \pm z\langle \phi|\langle \phi_{\alpha}|\mathbf{P}^{-1}\right)_{kl} = \left((\phi_{\alpha k}, \phi_{l}) \pm (\phi_{l}, \phi_{\alpha k})\right).
\tag{4.8.26}
\end{equation*}
用式 (4.8.26) 左端的两个函数作为新基，表示 $\mathbf{C}'_{\alpha}$ 变换为等价表示 $\mathbf{C}''_{\alpha}$：
\begin{equation*}
\mathbf{C}''_{\alpha}(l) = \begin{pmatrix} \mathbf{N}_{\alpha}(l) & \mathbf{0}\\ \mathbf{0} & \mathbf{N}_{\alpha}(l) \end{pmatrix},
\tag{4.8.27}
\end{equation*}
\begin{equation*}
\mathbf{C}''_{\alpha}(z) = \begin{pmatrix} \mathbf{P} & \mathbf{0}\\ \mathbf{0} & -\mathbf{P} \end{pmatrix}.
\tag{4.8.28}
\end{equation*}
于是，对称化基 $(\phi_{\alpha k}, \phi_{l}) + (\phi_{l}, \phi_{\alpha k})$ 给出 $\mathbf{M}_{\alpha}$ 的一个表示 $\mathbf{N}^{+}_{\alpha}$：
\begin{equation*}
\left.
\begin{array}{l}
\mathbf{N}^{+}_{\alpha}(l) = \mathbf{N}_{\alpha}(l)\\
\mathbf{N}^{+}_{\alpha}(z) = \mathbf{P}
\end{array}
\right\}
\tag{4.8.29}
\end{equation*}
而反对称化基 $(\phi_{\alpha k}, \phi_{l}) - (\phi_{l}, \phi_{\alpha k})$ 给出 $\mathbf{M}_{\alpha}$ 的一个表示 $\mathbf{N}^{-}_{\alpha}$：
\begin{equation*}
\left.
\begin{array}{l}
\mathbf{N}^{-}_{\alpha}(l) = \mathbf{N}_{\alpha}(l)\\
\mathbf{N}^{-}_{\alpha}(z) = -\mathbf{P}
\end{array}
\right\}.
\tag{4.8.30}
\end{equation*}
此外，当表示 $\mathbf{N}^{+}_{\alpha}$ 与 $\mathbf{N}^{-}_{\alpha}$ 从 $\mathbf{M}_{\alpha}$ 诱导到 $\mathbf{G}$ 时，它们显然分别定义在空间 $V^{+}_{\alpha}$ 与 $V^{-}_{\alpha}$ 上，其中 $V_{\alpha} = (V^{+}_{\alpha} \oplus V^{-}_{\alpha})$ 是 $V_{\alpha}$ 分解为其对称化与反对称化部分的直和。于是 $[(\mathbf{D} \uparrow \mathbf{G}) \boxtimes (\mathbf{D} \uparrow \mathbf{G})]$ 含有 $\mathbf{N}^{+}_{\alpha} \uparrow \mathbf{G}$，$\{(\mathbf{D} \uparrow \mathbf{G}) \boxtimes (\mathbf{D} \uparrow \mathbf{G})\}$ 含有 $\mathbf{N}^{-}_{\alpha} \uparrow \mathbf{G}$，并且如预期的那样这两个表示维数相等。

另外，由式 (4.8.22)、(4.8.29) 与 (4.8.30)，$\mathbf{N}^{+}_{\alpha}$ 与 $\mathbf{N}^{-}_{\alpha}$ 的特征标容易计算，对一切 $l \in \mathbf{L}_{\alpha}$ 为
\begin{equation*}
\chi_{N^{+}_{\alpha}}(l) = \chi_{D}(d^{-1}_{\alpha}ld_{\alpha})\chi_{D}(l)
\tag{4.8.31}
\end{equation*}
\begin{equation*}
\chi_{N^{\pm}_{\alpha}}(zl) = \pm\chi_{D}(zlzl).
\tag{4.8.32}
\end{equation*}
例如
\begin{equation*}
\begin{array}{l}
\sum_{i, j}\mathbf{N}^{-}_{\alpha}(zl)_{ij, ij} = \sum_{i, j, k, l}\mathbf{P}_{ij, kl}\mathbf{N}_{\alpha}(l)_{kl, ij}\\
= \sum_{i, j, k, l}\left\{\mathbf{D}(h') \otimes \mathbf{D}(h)\right\}_{ji, kl}\left\{\mathbf{D}(d^{-1}_{\alpha}ld_{\alpha}) \otimes \mathbf{D}(l)\right\}_{kl, ij}\\
= \sum_{i, j}\left\{\mathbf{D}(h'd^{-1}_{\alpha}ld_{\alpha}) \otimes \mathbf{D}(hl)\right\}_{ji, ij}\\
= \sum_{i, j}\mathbf{D}(h'd^{-1}_{\alpha}ld_{\alpha})_{ji}\mathbf{D}(hl)_{ij}\\
= \sum_{j}\mathbf{D}(h'd^{-1}_{\alpha}ld_{\alpha}hl)_{jj}\\
= \chi_{D}(zlzl).
\end{array}
\tag{4.8.33}
\end{equation*}
分解至此完成。把各种不同情形汇总，得到如下定理。

\noindent\textbf{定理 4.8.1}\quad 设 $\mathbf{D}$ 是有限群 $\mathbf{G}$ 的子群 $\mathbf{H}$ 的表示。设 $\mathbf{A}$ 是除 $\mathbf{H}$ 自身以外全部自逆双陪集 $\mathbf{H}d_{\alpha}\mathbf{H} = \mathbf{H}d^{-1}_{\alpha}\mathbf{H}$ 的集合。设 $\mathbf{B}$ 是全部形如 $\mathbf{H}d_{\beta}\mathbf{H} \cup \mathbf{H}d^{-1}_{\beta}\mathbf{H}$ 的集合的集合，其中 $\mathbf{H}d_{\beta}\mathbf{H}$ 是非自逆双陪集。对每个 $\beta \in \mathbf{B}$，令 $\mathbf{N}_{\beta}(l)$ 是 $\mathbf{L}_{\beta} = \mathbf{H} \cap d_{\beta}\mathbf{H}d^{-1}_{\beta}$ 的表示 $\mathbf{D}(d^{-1}_{\beta}ld_{\beta}) \otimes \mathbf{D}(l)$。对每个 $\alpha \in \mathbf{A}$，$d_{\alpha}\mathbf{H} \cap \mathbf{H}d^{-1}_{\alpha}$ 非空。设 $z$ 是它的任一成员，比如 $z = d_{\alpha}h = h'd^{-1}_{\alpha}$。设 $\mathbf{L}_{\alpha} = \mathbf{H} \cap d_{\alpha}\mathbf{H}d^{-1}_{\alpha}$，并设 $\mathbf{M}_{\alpha}$ 是由 $\mathbf{L}_{\alpha}$ 与 $z$ 生成的子群。则 $\mathbf{L}_{\alpha}$ 是 $\mathbf{M}_{\alpha}$ 的指数为 2 的不变子群。设 $\mathbf{N}_{\alpha}(l)$ 是 $\mathbf{L}_{\alpha}$ 的表示 $\mathbf{D}(d^{-1}_{\alpha}ld_{\alpha}) \otimes \mathbf{D}(l)$，并设 $\mathbf{P}$ 是由 $\mathbf{P}_{ab, cd} = \mathbf{D}(h')_{bc}\mathbf{D}(h)_{ad}$ 给出的矩阵。则存在 $\mathbf{N}_{\alpha}$ 到群 $\mathbf{M}_{\alpha}$ 上的扩充 $\mathbf{N}^{+}_{\alpha}$ 与 $\mathbf{N}^{-}_{\alpha}$，使得 $\mathbf{N}^{+}_{\alpha}(z) = \mathbf{P}$、$\mathbf{N}^{-}_{\alpha}(z) = -\mathbf{P}$。最后我们有
\begin{equation*}
\left[(\mathbf{D} \uparrow \mathbf{G}) \boxtimes (\mathbf{D} \uparrow \mathbf{G})\right] \equiv [\mathbf{D} \boxtimes \mathbf{D}] \uparrow \mathbf{G} + \sum_{\alpha \in \mathbf{A}}\mathbf{N}^{+}_{\alpha} \uparrow \mathbf{G} + \sum_{\beta \in \mathbf{B}}\mathbf{N}_{\beta} \uparrow \mathbf{G}
\tag{4.8.34}
\end{equation*}
\begin{equation*}
\left\{(\mathbf{D} \uparrow \mathbf{G}) \boxtimes (\mathbf{D} \uparrow \mathbf{G})\right\} \equiv \{\mathbf{D} \boxtimes \mathbf{D}\} \uparrow \mathbf{G} + \sum_{\alpha \in \mathbf{A}}\mathbf{N}^{-}_{\alpha} \uparrow \mathbf{G} + \sum_{\beta \in \mathbf{B}}\mathbf{N}_{\beta} \uparrow \mathbf{G}.
\tag{4.8.35}
\end{equation*}

作为例子，把上述定理应用于空间群 $\mathbf{G} = P23\ (T^{1})$，即 4.7 节的例子。现在的 $\mathbf{H}$ 是小群 $\mathbf{G}^{M}$。双陪集分解 $\mathbf{G} = \sum_{\lambda}\mathbf{G}^{M}d_{\lambda}\mathbf{G}^{M}$ 含三项，$d_{\lambda} = \{E \mid \mathbf{0}\}$、$\{C^{+}_{31} \mid \mathbf{0}\}$ 与 $\{C^{-}_{31} \mid \mathbf{0}\}$。于是集合 $\mathbf{A}$ 为空，集合 $\mathbf{B}$ 含单独一项 $\mathbf{G}^{M}\{C^{+}_{31} \mid \mathbf{0}\}\mathbf{G}^{M} \cup \mathbf{G}^{M}\{C^{-}_{31} \mid \mathbf{0}\}\mathbf{G}^{M}$。取 $D = A_{1}$，则在展开 (4.8.4) 中，$d_{\lambda} = d_{\varepsilon} = \{E \mid \mathbf{0}\}$ 的项诱导到 $\mathbf{G}$ 得到三个表示 $\Gamma A$、$\Gamma{}^{1}E$ 与 $\Gamma{}^{2}E$；$d_{\lambda} = \{C^{+}_{31} \mid \mathbf{0}\}$ 的项给出 $MA_{1}$；同样 $d_{\lambda} = \{C^{-}_{31} \mid \mathbf{0}\}$ 的项也给出 $MA_{1}$（这说明来自互逆双陪集的表示是等价的）——这些结果可用表 4.5 检验。又由于 $A_{1}$ 是一维的，$\{A_{1} \boxtimes A_{1}\}$ 为空。现在可以写下展开 (4.8.34) 与 (4.8.35)：
\begin{equation*}
\left[MA_{1} \boxtimes MA_{1}\right] \equiv \Gamma A + \Gamma{}^{1}E + \Gamma{}^{2}E + MA_{1},
\tag{4.8.36}
\end{equation*}
\begin{equation*}
\left\{MA_{1} \boxtimes MA_{1}\right\} \equiv MA_{1}.
\tag{4.8.37}
\end{equation*}
特别值得注意的是：当 Kronecker 平方的分解已知时，实践中为得到对称化与反对称化 Kronecker 平方的分解所需做的额外工作是何等之少；不过必须承认，上述情形非常有利——没有自逆双陪集，而且表示 $A_{1}$ 是一维的，使得 $\{A_{1} \boxtimes A_{1}\}$ 为零。
